\begin{tikzpicture}[
  code/.style={
    draw,
    rounded corners,
    fill=gray!8,
    font=\ttfamily\small,
    align=left,
    inner sep=8pt
  },
  note/.style={
    draw,
    rounded corners,
    fill=blue!6,
    font=\small,
    align=left,
    text width=4.4cm,
    inner sep=6pt
  },
  arrow/.style={->, thick}
]

\node[code] (theorem) {
theorem every\_element\_is\_one\_or\_successor\\
\quad (ps : PeanoSystem) :\\
\quad $\forall$ element : ps.carrier,\\
\quad\quad element = ps.one $\lor$\\
\quad\quad $\exists$ predecessor : ps.carrier,\\
\quad\quad\quad ps.successor predecessor = element := by\\[2mm]
\quad apply induction\_principle ps\\[1mm]
\quad $\cdot$ left\\
\quad\quad rfl\\[1mm]
\quad $\cdot$ intro element induction\_hypothesis\\
\quad\quad right\\
\quad\quad exact $\langle$element, rfl$\rangle$
};

\node[note, right=2.4cm of theorem.north east, anchor=north west] (ps) {
\textbf{\texttt{ps : PeanoSystem}}\\
The theorem is not about one fixed natural number system. It is about an arbitrary Peano system:
\[
(P,S,1).
\]
};

\node[note, right=2.4cm of theorem.east, anchor=west] (goal) {
\textbf{The target predicate}\\
Lean infers the predicate:
\[
P(x) := x=1 \;\lor\; \exists u,\; S(u)=x.
\]
This is the Lean-native version of defining the subset \(D\).
};

\node[note, right=2.4cm of theorem.south east, anchor=south west] (base) {
\textbf{Base case}\\
For \(x=1\), prove:
\[
1=1 \;\lor\; \exists u,\; S(u)=1.
\]
Choose the left side using \texttt{left}; prove \(1=1\) by \texttt{rfl}.
};

\node[note, below=1.3cm of theorem] (step) {
\textbf{Successor step}\\
To prove the statement for \(S(x)\), choose the right side:
\[
\exists u,\; S(u)=S(x).
\]
Take \(u=x\). Then \(S(x)=S(x)\), proved by \texttt{rfl}.
};

\node[note, left=2.4cm of theorem.west, anchor=east] (ih) {
\textbf{Induction hypothesis}\\
The variable \texttt{induction\_hypothesis} says the theorem already holds for \texttt{element}.\\
In this proof, it is not actually needed, because \(S(\texttt{element})\) is visibly a successor.
};

\draw[arrow] (ps.west) -- ++(-0.7,0) |- ([yshift=2.1cm]theorem.east);
\draw[arrow] (goal.west) -- ++(-0.7,0) |- ([yshift=0.6cm]theorem.east);
\draw[arrow] (base.west) -- ++(-0.7,0) |- ([yshift=-1.0cm]theorem.east);
\draw[arrow] (step.north) -- ([yshift=-1.7cm]theorem.south);
\draw[arrow] (ih.east) -- ++(0.7,0) |- ([yshift=-1.25cm]theorem.west);

\end{tikzpicture}